% !TeX root = ../navier-stokes-manuscript.tex
% ============================================================
% 2.1 Vortex Stretching
% ============================================================

A narrow material vortex tube stretched along its axis squeezes across its cross-section, drives its rotating fluid inward, and makes the interior spin faster.
That spin-up is the positive stretching contribution; whether the rotation strengthens overall is decided by the complete signed balance with viscosity.
The tube moves inside one velocity field whose local acceleration, nonlinear transport, pressure, and viscous diffusion obey
\[
  \underbrace{\partial_t u}_{\text{local acceleration}}
  +
  \underbrace{(u\cdot\nabla)u}_{\text{nonlinear transport}}
  =
  \underbrace{-\nabla p}_{\text{pressure}}
  +
  \underbrace{\nu\Delta u}_{\text{viscous diffusion}}.
\]
Its narrowing and spin-up now have to be read within this full evolution, where transport carries the tube and the signed viscous contribution can oppose the stretching.

\subsection{Vortex Stretching}\label{sub:ThePerfectlyStretchingVortex}

Let the tube be locally axisymmetric, with vorticity aligned with its axis, and suppose the strain is locally uniform across it.
If \(e\) is that axis and \(L(t)\) is the tube's length, its aligned extension is described by
\[
  S:=\frac12\bigl(\nabla u+(\nabla u)^{\mathsf T}\bigr),
  \qquad
  Se=\sigma(t)e,
  \qquad
  \sigma(t)>0,
  \qquad
  \frac{\dot L(t)}{L(t)}
  =
  e\cdot Se
  =
  \sigma(t).
\]
Incompressibility turns that extension into transverse contraction.
Its material volume \(\mathcal V\) stays fixed, so define its effective transverse area by \(A(t):=\mathcal V/L(t)\) and its effective radius by \(r_{\mathrm{eff}}(t):=\sqrt{A(t)/\pi}\).
Then
\[
  \nabla\cdot u=0,
  \qquad
  \frac{d}{dt}(A L)=0,
  \qquad
  \frac{\dot A}{A}=-\sigma(t),
  \qquad
  \frac{\dot r_{\mathrm{eff}}}{r_{\mathrm{eff}}}
  =
  -\frac{\sigma(t)}{2}.
\]
Thus the same positive axial strain that lengthens the tube draws its rotating fluid toward the axis.
With \(\omega:=\nabla\times u\), the vorticity carried along that axis satisfies
\[
  \frac{D\omega}{Dt}
  =
  S\omega+\nu\Delta\omega,
  \qquad
  \frac{D}{Dt}
  :=
  \partial_t+u\cdot\nabla,
\]
and under the alignment \(\omega=|\omega|e\),
\[
  S\omega
  =
  \sigma(t)\omega,
  \qquad
  \frac12\frac{D|\omega|^2}{Dt}
  =
  \sigma(t)|\omega|^2
  +
  \nu\,\omega\cdot\Delta\omega.
\]
The first term records the stretching spin-up, while the full right-hand side decides its net result.
On an interval where that side is positive, the rotation strengthens as the effective radius contracts.
The energy identity proved in \Cref{prop:ClassicalKineticEnergyIdentity} and the antisymmetric part of \(\nabla u\) give
\[
  \norm{u(t)}_2
  \leq
  \norm{u_0}_2
  <\infty,
  \qquad
  \left|\frac12\bigl(\nabla u-(\nabla u)^{\mathsf T}\bigr)\right|_{\mathrm F}
  =
  \frac{|\omega|}{\sqrt2}
  \leq
  |\nabla u|_{\mathrm F}.
\]
Finite kinetic energy controls the global size of the velocity, but it does not close the possibility of concentrated local growth in its rotational gradient.

\divider{\NavierStokes{} Scaling}

To test whether a finer scale changes the balance, fix \(t_0<T_*\) and call the tube's effective radius at that time \(\ell:=r_{\mathrm{eff}}(t_0)\).
For \(\lambda>1\), the exact \NavierStokes{} comparison at width \(\lambda^{-1}\ell\) is the rescaled solution
\[
  u_\lambda(x,t)
  :=
  \lambda u(\lambda x,\lambda^2t),
  \qquad
  p_\lambda(x,t)
  :=
  \lambda^2p(\lambda x,\lambda^2t).
\]
At time \(\lambda^{-2}t_0\), the corresponding segment in \(u_\lambda\) has effective radius \(\lambda^{-1}\ell\).
This is a complete solution at another scale, rather than a later state of the original material tube.
Its four momentum terms transform together:
\[
\begin{aligned}
  \partial_tu_\lambda(x,t)
  &=
  \lambda^3(\partial_tu)(\lambda x,\lambda^2t),
  \\
  (u_\lambda\cdot\nabla)u_\lambda(x,t)
  &=
  \lambda^3\bigl((u\cdot\nabla)u\bigr)(\lambda x,\lambda^2t),
  \\
  \nabla p_\lambda(x,t)
  &=
  \lambda^3(\nabla p)(\lambda x,\lambda^2t),
  \\
  \nu\Delta u_\lambda(x,t)
  &=
  \lambda^3\nu(\Delta u)(\lambda x,\lambda^2t).
\end{aligned}
\]
Nonlinear transport and viscosity both gain the same \(\lambda^3\) factor, as do local acceleration and pressure.
Passing to a finer scale gives viscosity no relative advantage by scaling alone.

\divider{Critical Height}

Although the equation keeps this balance, its homogeneous Sobolev norms change at different rates.
Let \(\Lambda:=(-\Delta)^{1/2}\).
Since
\[
  \widehat{u_\lambda}(\xi,t)
  =
  \lambda^{-2}\widehat u(\lambda^{-1}\xi,\lambda^2t),
\]
a change of variables gives
\[
  \norm{u_\lambda(t)}_{\dot H^s}^2
  =
  \lambda^{2s-1}
  \norm{u(\lambda^2t)}_{\dot H^s}^2.
\]
Within this displayed homogeneous Sobolev family, kinetic energy at \(s=0\) loses one power of \(\lambda\), while the first-derivative level at \(s=1\) gains one.
The exponent vanishes at \(s=\tfrac12\).
Set
\[
  H(t)
  :=
  \frac12\norm{u(t)}_{\dot H^{1/2}}^2
  =
  \frac12\norm{\Lambda^{1/2}u(t)}_2^2.
\]
Writing \(H_\lambda\) for the same quantity evaluated on \(u_\lambda\),
\[
  \norm{u_\lambda(t)}_2^2
  =
  \lambda^{-1}\norm{u(\lambda^2t)}_2^2,
  \qquad
  H_\lambda(t)=H(\lambda^2t),
  \qquad
  \norm{\nabla u_\lambda(t)}_2^2
  =
  \lambda\norm{\nabla u(\lambda^2t)}_2^2.
\]
The rescaled velocity has less kinetic energy and a larger first-derivative norm, while its half-derivative height is scale invariant.
Scaling identifies this critical height, but does not determine whether it rises or falls in time.

\divider{Nonlinear Production versus Viscous Dissipation}

At critical height, the nonlinear contribution is the signed production
\[
  \mathcal P_{1/2}[u](t)
  :=
  -\operatorname{Re}
  \pair
    {\Lambda^{1/2}u(t)}
    {\Lambda^{1/2}\bigl((u(t)\cdot\nabla)u(t)\bigr)}_{L^2}.
\]
The pressure pairing vanishes because \(\nabla\cdot u=0\), and the Laplacian contributes \(-\nu\norm{\Lambda^{3/2}u}_2^2\).
The critical height obeys
\[
  H'(t)
  +
  \nu\norm{\Lambda^{3/2}u(t)}_2^2
  =
  \mathcal P_{1/2}[u](t).
\]
Its sign is fixed by the signed production after the positive viscous dissipation is subtracted: \(H\) rises exactly when production exceeds dissipation.
Under the same rescaling,
\[
  \mathcal P_{1/2}[u_\lambda](t)
  =
  \lambda^2\mathcal P_{1/2}[u](\lambda^2t),
  \qquad
  \nu\norm{\Lambda^{3/2}u_\lambda(t)}_2^2
  =
  \lambda^2\nu\norm{\Lambda^{3/2}u(\lambda^2t)}_2^2.
\]
Production and dissipation remain matched under scale, both acquiring the inverse of the time-contraction factor.
Their relative scaling leaves open the physical time on which a localized structure changes.

\divider{The Heat Clock}

For a scale-localized velocity structure, the heat equation \(v_t=\nu\Delta v\) isolates the linear viscous response:
\[
  \widehat v(\xi,t+\delta t)
  =
  e^{-\nu|\xi|^2\delta t}\widehat v(\xi,t).
\]
A structure of width \(r\) occupies frequencies \(|\xi|\simeq r^{-1}\).
The heat multiplier changes it by order one when
\[
  \nu|\xi|^2\delta t\simeq1,
\]
so its linear viscous comparison time is
\[
  \tau_\nu(r)
  :=
  \frac{r^2}{\nu}.
\]
This time measures diffusion acting on the scale-localized velocity structure, not the lifetime of the opening material tube.
Halving the localization width quarters the comparison time.
More generally,
\[
  \tau_\nu(\lambda^{-1}r)
  =
  \lambda^{-2}\tau_\nu(r).
\]
The same quadratic contraction appears in the full \NavierStokes{} rescaling.

\divider{The Zeno Cascade}

Now take dyadic localization widths.
For
\[
  r_n:=2^{-n}r_0,
  \qquad
  \tau_n:=\frac{r_n^2}{\nu},
\]
the corresponding heat times satisfy
\[
  \tau_n
  =
  \frac{r_0^2}{\nu}4^{-n},
  \qquad
  \sum_{n=0}^{\infty}\tau_n
  =
  \frac{4r_0^2}{3\nu}
  <
  \infty.
\]
This finite arithmetic becomes a candidate fluid history if one \NavierStokes{} velocity field exhibits scale-localized structures at widths
\[
  r_0>r_1>r_2>\cdots\longrightarrow0
\]
and at sampled times
\[
  t_0<t_1<t_2<\cdots\uparrow T_*<\infty,
  \qquad
  t_{n+1}-t_n\asymp\tau_n,
\]
with comparison constants independent of \(n\).
Under that one-field condition, the remaining time satisfies
\[
  T_*-t_n
  \asymp
  \sum_{k=n}^{\infty}\tau_k
  =
  \frac{4r_n^2}{3\nu}.
\]
At width \(r_n\), both the next sampled gap and the total time left are comparable to \(r_n^2/\nu\).
The condition does not establish that such a cascade occurs, but it shows the temporal demand a single field would have to meet: its scale-localized structures would change across intervals shrinking to zero as \(t\uparrow T_*\).
Those collapsing intervals leave the question that comes next: how do that same field's first and successive Eulerian time derivatives behave at one fixed spatial point?
